\documentclass[10pt,a4paper]{amsart}
\usepackage[margin=19mm]{geometry}
\usepackage{amsmath,amssymb,mathtools}
\usepackage[scr=rsfs,cal=euler]{mathalfa}
\usepackage{xfrac}
\usepackage[unicode,hidelinks,bookmarks=false]{hyperref}
\newcommand{\ie}{i.e.\ }
\newcommand{\cf}{cf.\ }
\newcommand{\resp}{respectively\ }
\newcommand{\Subsection}{\subsection}
\providecommand{\nobreakdash}{\nobreak\hbox{-}}
\setlength{\emergencystretch}{2em}
\multlinegap0pt
\usepackage{bm}
\usepackage{bbm}
\usepackage{dsfont}
\usepackage{xspace}
\usepackage{cancel}
\usepackage{mathtools}
\usepackage{amsthm}
\makeatletter
\@ifundefined{theorem}{\newtheorem{theorem}{Theorem}}{}
\@ifundefined{claim}{\newtheorem{claim}[theorem]{Claim}}{}
\makeatother
\DeclareMathOperator{\supp}{supp}
\def\N{\mathbb{N}}
\def\vol{\mathrm{Vol}}
\newcommand{\Cone}{\mathrm{Cone}}
\newcommand{\SL}{\mathrm{SL}}
\newcommand{\SO}{{\mathrm{SO}}}
\title[Quantitative Oppenheim Conjecture for Random Quadratic Forms]{Quantitative Oppenheim Conjecture for Random Quadratic Forms and Optimal Variance Bounds in Function Fields}
\author{Jiyoung Han, Noy Soffer Aranov}
\date{Source: arXiv:2606.16699v1 (CC BY 4.0)}
\begin{document}
\maketitle

\noindent\emph{Source and licence.} Quantitative Oppenheim Conjecture for Random Quadratic Forms and Optimal Variance Bounds in Function Fields. Version arXiv:2606.16699v1 (CC BY 4.0), distributed under the Creative Commons Attribution 4.0 International licence, as recorded in the arXiv licence metadata for this exact version and shown on the abstract page https://arxiv.org/abs/2606.16699v1. Full rights notice: licenses/arxiv-2606.16699-CC-BY-4.0.txt.\par\smallskip

\noindent\textit{Editorial scope.} Counted unit: the Theorem whose source label is \texttt{thm: Witt} in the article (source0.tex lines 942--945, source label \texttt{thm: Witt}), delivered together with the article's own complete original proof of it (source0.tex lines 946--1110, one \texttt{proof} environment of 165 lines). No mathematics is written, repaired or re-derived: statement and proof are the article's own text line for line. Required context retained uncounted: lem:UM (lines 143--146); lem:UM= (lines 143--146); volume estimate (lines 905--930); standard quadratic form (lines 907--909); eqn: Sobolev (lines 912--914); eqn:smoothVolEst (lines 917--921); thm:VolR=0 (lines 936--940). Every definition, display and label of those lines is retained unchanged. Omitted from this package and not redistributed: 1--142, 147--904, 910--911, 915--916, 922--935, 941--941, 1111--1384. Nothing inside a retained excerpt is omitted or altered: every retained body is delivered exactly as the article's lines are written, so no omission receipt of any kind accompanies this package. Citations: the retained excerpts carry no citation command at all, so no bibliography entry of the article is transcribed and no reference is redistributed. Reference adaptation: none was needed and none was made; every reference inside the delivered unit resolves inside it, so no unresolved reference or citation is produced. Printed slips kept literal: no printed word, symbol, hypothesis or number of the counted statement or of its proof was altered. Layout and class: the article's own class is \texttt{amsart} with packages including amssymb, amsmath, amsthm, fontenc, enumerate, xy, fullpage, comment, array, longtable, stmaryrd, mathrsfs, xcolor, mathtools; the wrapper is the lane's pinned \texttt{amsart} editorial wrapper at 10pt on A4, which restates the theorem-like environments the retained text needs and adds the article's own macro definitions that the retained text needs (\texttt{\textbackslash Cone}, \texttt{\textbackslash N}, \texttt{\textbackslash SL}, \texttt{\textbackslash SO}, \texttt{\textbackslash supp}, \texttt{\textbackslash vol}), with extra wrapper packages (bm, bbm, dsfont, xspace, cancel, mathtools). The delivered numbering is the unit's own. Assets: no retained span contains a figure environment, a graphics-inclusion command, a PDF-page inclusion, a table or a TikZ picture; no image file is copied into this package, the package declares no asset dependency and no image file is redistributed with it. Licence: version arXiv:2606.16699v1 of this article is distributed under the Creative Commons Attribution 4.0 International licence, as recorded in the arXiv licence metadata for this exact version and shown on the abstract page https://arxiv.org/abs/2606.16699v1. A full rights notice is delivered at licenses/arxiv-2606.16699-CC-BY-4.0.txt. Characters: every retained excerpt is pure ASCII, so the delivered engine, XeLaTeX with its default Latin Modern text font, reports no missing character.
    \begin{enumerate}
        \item \label{lem:UM} $\Vert \mathbf{u}+\mathbf{v}\Vert\leq \max\{\Vert \mathbf{u}\Vert,\Vert\mathbf{v}\Vert\}$;
        \item \label{lem:UM=} If $\Vert \mathbf{u}\Vert\neq \Vert \mathbf{v}\Vert$, then $\Vert \mathbf{u}+\mathbf{v}\Vert=\max\{\Vert \mathbf{u}\Vert,\Vert \mathbf{v}\Vert\}$.
    \end{enumerate}

\begin{theorem}\label{volume estimate}
Let $Q_0$ be a non-degenerate isotropic quadratic form given as
\begin{equation}\label{standard quadratic form}
Q_0(v_1, \ldots, v_n)=2v_1v_n+Q'(v_2, \ldots, v_{n-1}),
\end{equation}
where $Q'$ is some quadratic form of $(n-2)$ variables. Assume that the absolute values of all the coefficients of $Q_0$ are bounded by $q^N$ for some $N\ge 0$.
Let $h\in C_c^{\infty}(\mathcal K_\infty^n)$ be such that $\supp h \subseteq x^{R}\mathcal O^n\setminus x^{r-1}\mathcal O^n$, where $r\le R\in \mathbb Z$, and
\begin{equation}\label{eqn: Sobolev}
h(\mathbf v + x^{-K}\mathbf{u})=h(\mathbf v), \;\forall \mathbf v\in \mathcal K_\infty^n,\forall \mathbf{u}\in \mathcal{O}^n
\end{equation}
for some positive integer $K\in \mathbb N$. We may assume that $K\ge -r$.
Let $I=a+x^b\mathcal O\subseteq \mathcal K_\infty$ and let $B=\max\{\log_q|a|, b\}$. Then, for any $t> \frac{\max\left\{-2r+B,N+B-b,-3r+B+K\right\}} 2$, we have
\begin{equation}\begin{split}
\label{eqn:smoothVolEst}
\int_{\mathcal K_\infty^n} h(x^{-t}\mathbf v) \chi^{}_{I}(Q_0(\mathbf v)) d\mathbf v
&=J_{h,Q_0}q^{b}q^{(n-2)t},\\
\end{split}\end{equation}
where $J_{h,Q_0}>0$ is a constant depending only on $h$ and $Q_0$.

As a consequence, for any quadratic form $Q\in \{Q_0^g: g\in \SL_n(\mathcal K_\infty)\}$ there exists a constant $c_Q>0$, such that for any $I$ as above and for any $T>\frac{1}{2}\max\left\{0,B,N+B-b,B+K\right\}$, where $B,N,K$ are as above, we have
\begin{equation}
\label{eqn:VolEst}
\vol\left\{\mathbf v \in \mathcal K_\infty^n: \|\mathbf v \|\leq q^T\;\text{and}\; Q(\mathbf v)\in I\right\}
=c_Q \vol(I)q^{(n-2)T}+O_{Q,I}\left(1\right).
\end{equation}
\end{theorem}

\begin{theorem}
\label{thm:VolR=0}
    Let $Q_0$ be as in Theorem \ref{volume estimate}. Let $h\in C_c^{\infty}(\mathcal{K}_{\infty}^n)$ be such that $\operatorname{supp}(h)\subseteq \mathcal{O}^n\setminus \mathfrak{m}^n$, and assume that there exists some $K\in \mathbb{N}$, such that \eqref{eqn: Sobolev} holds. Let $I=a+x^{b}\mathcal{O}\subseteq \mathcal{K}_{\infty}$, and let $B=\max\{\log_q\vert a\vert,b\}$. Then, there exists $J_{h,Q_0}$, such that for any $t>\frac{1}{2}\max\{B,N+B-b,B+K\}$, we have
    $$\int_{\mathcal{K}_{\infty}^n}h(x^{-t}\mathbf{v})\chi_I(Q_0(\mathbf{v}))d\mathbf{v}=J_{h,Q_0}q^bq^{(n-2)t}.$$
\end{theorem}

\begin{theorem}\label{thm: Witt}
    Let $\alpha\in \mathcal{K}_{\infty}$ and let $z\in\mathbb{Z}$. The group $\mathrm{K}=\operatorname{SO}(Q)\cap \operatorname{SL}_n(\mathcal{O}_{\infty})$ acts transitively on the set
    $$L(\alpha,q^z):=\left\{\mathbf{v}\in \mathcal{K}_{\infty}^n:Q(\mathbf{v})=\alpha,\Vert \mathbf{v}\Vert=q^z\right\}.$$
\end{theorem}

\begin{proof}[Proof of Theorem \ref{thm:VolR=0}]
Since smooth functions can be approximated with sums of finitely many characteristic functions, we may assume that $h$ is the characteristic function of a Borel set $A\subseteq \mathcal{O}^n\setminus \mathfrak{m}^n$ satisfying
$$A+x^{-K}\mathcal{O}^n=A.$$
By defining $x^{-t}\mathbf v=\mathbf w$, we have
\begin{equation}\label{eqn 1: volume estimate1}
\int_{\mathcal K_\infty^n} h(x^{-t}\mathbf v) \chi^{}_I(Q_0(\mathbf v)) d\mathbf v
=q^{nt}\int_{\mathcal K_\infty^n} h(\mathbf w) \chi^{}_I(x^{2t}Q_0(\mathbf w)) d\mathbf w.
\end{equation}
Set
\[
C_t(A,I)=\left\{\mathbf w\in \mathcal K_\infty^n: \mathbf w \in A\;\text{and}\;  x^{2t}Q_0(\mathbf w)\in I\right\},
\]
so that
\[
\eqref{eqn 1: volume estimate1}= q^{nt} \vol(C_t(A,I)).
\]
\begin{claim}
For every $t>\frac{1}{2}\max\{B+K, B+N\}$, we have
\[
C_t(A,I)=C'_t(A,I),
\]
where
\[
C'_t(A,I)=\left\{\mathbf u_1+ x^{-2t} \mathbf u_2: \begin{array}{c}
\mathbf u_1\in A\cap \Cone(Q_0);\\[0.05in]
\|\mathbf u_2\|\le q^B,\; 2Q_0(\mathbf u_1, \mathbf u_2)\in I
\end{array}\right\},
\]
where by abuse of notation, $Q_0$ denotes the bilinear form corresponding to the quadratic form $Q_0$, that is, $$Q_0(\mathbf v, \mathbf w)=\frac 1 2\left(Q_0(\mathbf v+\mathbf w)-Q_0(\mathbf v)-Q_0(\mathbf w)\right),$$
so that $Q_0(\mathbf{v},\mathbf{v})=Q_0(\mathbf{v})$.
\end{claim}
\begin{proof}
The fact that $C'_t(A,I) \subseteq C_t(A,I)$ follows from the following observations: Let $\mathbf{u}_1+x^{-2t}\mathbf{u}_2\in C_t'(A,I)$.
\begin{itemize}
\item Since $2t>B$ and $\|\mathbf u_2\|\le q^B$, we have 
\[\|x^{-2t} \mathbf u_2\|<1=\|\mathbf u_1\|.\] 
Thus, by \eqref{lem:UM=}, $\Vert \mathbf{u}_1+x^{-2t-z}\mathbf{u}_2\Vert=1$.
\vspace{1pt}
\item Since $2t> B+K$, we have $x^{-2t}\mathbf u_2 \in x^{-K}\mathcal O^n$. Thus, the fact that $\mathbf{u}_1\in A$ and \eqref{lem:UM=} imply that $\mathbf{u}_1+x^{-2t}\mathbf{u}_2\in A$. 
\vspace{1pt}
\item Since $2t\geq N+B-b$, we have 
$$x^{-2t}|Q_0(\mathbf u_2)|\leq q^{-2t+N}\Vert \mathbf{u}_2\Vert^2\le q^{-2t+N+2B}\le q^b.$$ Hence, by \eqref{lem:UM},
\begin{equation}
\begin{split}
\left| x^{2t}Q_0(\mathbf u_1 + x^{-2t}\mathbf u_2)-a\right|
=\left| 2Q_0(\mathbf u_1, \mathbf u_2)+x^{-2t}Q_0(\mathbf u_2)-a\right|\\
\leq \max\left\{\left| 2Q_0(\mathbf{u}_1,\mathbf{u}_2)-a\right|,\left|x^{-2t}Q_0(\mathbf{u}_2)\right|\right\}\leq q^b.
\end{split}
\end{equation}
\end{itemize}
Thus, $x^{2t}Q_0(\mathbf{u}_1+x^{-2t}\mathbf{u}_2)\in I$, so that $\mathbf{u}_1+x^{-2t}\mathbf{u}_2\in C_t(A,I)$. 
\vspace{5pt}

\noindent To prove that $C_t(A,I)\subseteq C_t'(A,I)$, note that by Theorem~\ref{thm: Witt}, the group $\mathrm{K}=\SO(Q_0)\cap \operatorname{SL}_n(\mathcal{O})$ acts transitively on the set $\left\{\mathbf{v}\in \mathcal{K}_{\infty}^n:Q_0(\mathbf{v})=Q_0(\mathbf{w}),\Vert \mathbf{v}\Vert=\Vert \mathbf{w}\Vert\right\}$. Moreover, by \eqref{standard quadratic form},
$$Q_0\left(\mathbf{e}_1+\frac{Q_0(\mathbf{w})}{2}\mathbf{e}_n\right)=2\cdot \frac{Q_0(\mathbf{w})}{2}=Q_0(\mathbf{w}).$$
In addition, note that $$\vert Q_0(\mathbf{w})\vert\leq \max\left\{q^{-2t}\vert a\vert,\left|Q_0(\mathbf{w})-\frac{a}{x^{2t}}\right|\right\}\leq \max\{q^{B-2t},q^{b-2t}\}<1.$$ Thus, by \eqref{lem:UM=}, $\Vert\mathbf{e}_1+Q_0(\mathbf{w})\mathbf{e}_n\Vert=\Vert \mathbf{e}_1\Vert$, and hence, $\|\mathbf u_1+x^{-2t}\mathbf u_2\|=\|\mathbf u_1\|=1$.
Thus, for every $\mathbf w\in C_t(A,I)$, there exists $k\in \mathrm{K}$ such that
\[
k^{-1}\mathbf w=\mathbf{e}_1+\frac{Q_0(\mathbf{w})}{2}\mathbf{e}_n=\mathbf e_1+ x^{-2t}\:\frac {x^{2t}Q_0(\mathbf w)} 2\mathbf e_n.
\]
Set $\mathbf u_1=k\mathbf e_1$ and $\mathbf u_2=\frac {x^{2t}Q_0(\mathbf w)} 2 k\mathbf e_n$. 


\begin{itemize} 
\item Since $Q_0(\mathbf w)\in x^{-2t}a+x^{-2t+b}\mathcal O\subseteq x^{-2t+B}\mathcal{O}$,
\[
\|\mathbf u_2\|=\left| \frac {x^{2t}Q_0(\mathbf w)}{2}\right|\cdot \Vert k\mathbf{e}_n\Vert\le q^{2t}q^{-2t+B}\cdot 1=q^B.
\]
\vspace{1pt}
\item Thus, for $2t>B+K$, since $\mathbf u_1=\mathbf w- x^{-2t}\mathbf u_2=k\mathbf{e}_1$ and $\mathbf w\in A$,
\[
\mathbf u_1 \in A \cap \Cone(Q_0).
\]
Clearly, $2Q_0(\mathbf u_1, \mathbf u_2)=Q_0(\mathbf w)\in I$.
\end{itemize}
Hence, $\mathbf{w}=\mathbf{u}_1+x^{-2t}\mathbf{u}_2\in C_t(A,I)$.
\vspace{5pt}

\noindent We now compute the volume of $C'_t(A,I)$. Note that the volume on $\mathcal K_\infty^n$ is the limit of the normalized counting measures of the image set of the projection $\mathcal K_\infty^n \rightarrow \mathcal K_\infty^n/ x^{-\ell}O^n$ for $\ell\in \N$, i.e.,

\begin{equation}
\begin{split}\label{eqn:Vol(C_t'(A,I)))}
&\vol\left(C'_t(A,I)\right)\\
&=\lim_{\ell\rightarrow \infty} \frac 1 {q^{n\ell}} \:
 \#\left\{[\mathbf u_1 + x^{-2t} \mathbf u_2]_\ell: 
 \begin{array}{c}
 \mathbf u_1\in A\cap \Cone(Q_0),\\
 {[\mathbf{u}_2]_{\ell}\in [x^B\mathcal{O}^n]_{\ell}},\; [2Q_0(\mathbf u_1, \mathbf u_2)]_{\ell}\in [I]_{\ell}\end{array}
 \right\}\\[0.1in]
 &=\lim_{\ell\rightarrow \infty} \frac 1 {q^{n\ell}}
\sum_{\scriptsize\begin{array}{r}
[\mathbf{u}_1]_{\ell}\in [\mathcal{O}^n\setminus x^{-1}\mathcal{O}^n]_{\ell}\\
\cap [A\cap \operatorname{Cone}(Q_0)]_{\ell}\end{array}}
\frac{\# 
 \left\{[x^{-2t}\mathbf u_2]_\ell\in [x^{B-2t}\mathcal{O}^n]_{\ell}: 2Q_0(\mathbf u_1, \mathbf u_2)\in I \right\}
}{\# 
    \left\{\begin{array}{l}[(\mathbf{v}_1,\mathbf{v}_2)]_{\ell}\in [A\cap \operatorname{Cone}(Q_0)]_{\ell}\times [x^B\mathcal{O}^n]_{\ell}: \\[0.05in]
\hspace{0.2in}\mathbf v_1+x^{-2t}\mathbf v_2\in C'_t(A,I)\;\text{s.t.}\\[0.05in]
\hspace{1.2in}\mathbf v_1+x^{-2t}\mathbf v_2=\mathbf u_1+x^{-2t}\mathbf u_2\end{array}\right\}
},
\end{split}\end{equation}
where $[\mathbf u]_\ell$ is the equivalence class of $\mathbf u\in \mathcal K_\infty^n $ modulo $x^{-\ell}$ and $[A]_{\ell}=\{[\mathbf v]_\ell: \mathbf u\in A\}$. We now compute the numerator and denominator of the last line of \eqref{eqn:Vol(C_t'(A,I)))} through the following two facts.

\vspace{0.1in}
\noindent i) For any $\mathbf u_1\in A\cap \Cone(Q_0)$ and $\ell>2t-b$, we have 
\[
\#\{[x^{-2t}\mathbf u_2]_\ell: 2Q_0( \mathbf u_1, \mathbf u_2)\in I\}=q^{(-2t+\ell)n+B(n-1)}\cdot \vol(I).
\] 
\begin{proof}
By Theorem \ref{thm: Witt}, $\mathrm{K}$ acts transitively on $\left\{\mathbf{v}:\Vert \mathbf{v}\Vert=1\right\}\cap \Cone(Q_0)$. Thus, there exists $k\in \mathrm{K}$, such that $\mathbf{u}_1=k\mathbf{e}_1$. Hence, by replacing $A$ with the set $kA$,  we can assume that $\mathbf u_1=\mathbf e_1$. 

Denote $\mathbf u_2=(u_1, \ldots, u_n)$. By \eqref{standard quadratic form},
\[
2Q_0(\mathbf e_1, \mathbf u_2)\in I
\quad\Leftrightarrow\quad
2u_n\in I,
\]
that is, $u_n\in \frac a 2  + x^{b}\mathcal O$. Thus 
\[\begin{split}
&\#\{[x^{-2t}\mathbf u_2]_\ell: [2Q_0( \mathbf u_1, \mathbf u_2)]_{\ell}\in [I]_{\ell}\}
=\#\left[x^{-2t}\left(x^B\mathcal O^{n-1}\times \left(\frac a 2 +x^{b}\mathcal O\right)\right)\right]_{\ell}\\
&\hspace{0.4in}=q^{(-2t+\ell)n}\cdot q^{B(n-1)} \cdot q^{b}=q^{(-2t+\ell)n}\cdot q^{B(n-1)} \cdot \:\vol(I).
\end{split}\]
\end{proof}

\vspace{0.1in}
\noindent ii) For any $\mathbf u_1,\mathbf v_1$ in $A\cap \Cone(Q_0)$ and $\mathbf{u}_2,\mathbf{v}_2\in x^B\mathcal{O}^n$, such that $\mathbf{u}_1+x^{-2t}\mathbf{u}_2\in C_t'(A,I)$, we have
\begin{equation}
\label{eqn:u_1=v_1 modx^-2t+B}
\mathbf u_1 + x^{-2t}\mathbf u_2=\mathbf v_1 + x^{-2t}\mathbf v_2\;
\quad\Leftrightarrow\quad
\mathbf u_1\equiv \mathbf v_1 \mod x^{-2t+B}.
\end{equation}
By Theorem \ref{thm: Witt}, by replacing $A$ with $kA$ for some $k\in \mathrm{K}$, we can assume that $\mathbf u_1=\mathbf e_1$ and compute the cardinality of the equivalence class of $\mathbf e_1$ modulo $x^{-2t+B}$ which intersects $\Cone(Q_0)$. Thus, by \eqref{eqn:u_1=v_1 modx^-2t+B},
\begin{equation}\label{eqn 2: volume estimate1}\begin{split}
&\# \left\{\left([\mathbf v_1]_{\ell},[x^{-2t}\mathbf v_2]_\ell\right): \begin{array}{c}
[\mathbf v_1+x^{-2t}\mathbf v_2]_{\ell}\in [C'_t(A,I)]_{\ell}\;\text{s.t.}\\[0.05in]
[\mathbf v_1+x^{-2t}\mathbf v_2]_{\ell}=[\mathbf u_1+x^{-2t}\mathbf u_2]_{\ell}\\
\text{ and } [\mathbf{v}_1]_{\ell}\in [\operatorname{Cone}(Q_0)]_{\ell}\end{array}\right\}\\[0.05in]
&\hspace{1in}=\#\left\{[ \mathbf e_1+x^{-2t+B}\mathbf z]_\ell: \|\mathbf z\|\le 1 \right\}\cap \Cone(Q_0).
\end{split}\end{equation}

By expanding $Q_0( \mathbf e_1+x^{-2t+B}\mathbf z)=0$, the vector $\mathbf z=(z_1, \ldots, z_n)$ must satisfy the following equation 
\[
x^{-2t+B}z_n=\frac {-Q'(x^{-2t+B}z_2, \ldots,x^{-2t+B}z_{n-1})}{2(1+x^{-2t+B}z_1)}.
\]
Hence 
\[
\eqref{eqn 2: volume estimate1}=q^{(-2t+B+\ell)(n-1)}.
\]

As a consequence, 
\[\begin{split}
\eqref{eqn 1: volume estimate1}
&=q^{nt}\vol\left(C'_t(A,I)\right)\\
&=\lim_{\ell\rightarrow \infty}\frac {q^{nt}}{q^{n\ell}}\cdot \frac{\#[A\cap \Cone(Q_0)]_\ell\times q^{(-2t+\ell)n+B(n-1)}\cdot \vol(I)}{q^{(-2t+B+\ell)(n-1)}}\\
&=q^{(n-2)t}\: \vol(I)\lim_{\ell\rightarrow \infty} \frac {\#[A\cap \Cone(Q_0)]_\ell}{q^{(n-1)\ell}}
\end{split}\]
\end{proof}
The asymptotic formula \eqref{eqn:smoothVolEst} follows if we define
\[
J_{h,Q_0}=\lim_{\ell\rightarrow \infty} \frac {\#[A\cap \Cone(Q_0)]_\ell}{q^{(n-1)\ell}},
\]
where $J_{h,Q_0}$ is well-defined for general $h\in C_c(\mathcal O^n_\infty\setminus x^{-1}\mathcal O^n)$.
\end{proof}

\end{document}
